Agradeço a Deus, pela fé que me sustentou ao longo do curso e deste trabalho.

À minha família, pelo apoio e pela paciência durante os anos de graduação.

Ao meu orientador, Prof. Dr. Rafael Silva Guimarães, pela orientação e pelo acompanhamento ao longo deste trabalho. Ao Prof. Dr. Everson Scherrer Borges, que orientou a sua primeira etapa.

À Marcelli da Fonseca Carvalho, que aceitou julgar, uma a uma e sem saber de onde vinham, as marcações da auditoria cega, sem obrigação nenhuma de fazê-lo. O trabalho se propôs a medir o critério de marcação contra o julgamento de alguém que não participou da calibração, e esse olhar foi o dela. Obrigado pelo tempo e pela paciência com centenas de recortes de rocha.

Aos colegas de pesquisa com quem construí o DeepStoneAI e o FracStoneAI, trabalhos dos quais o ARIA é continuação: Matheus Gimenes de Souza, Arthur Cesar Figueiredo da Silva, Andrey Magalhães Silva, Pedro Brito Moreira, Igor Henrique Beloti Pizetta e Juliano Tessinari Zagôto. A João Victor Costa Araujo, por ter tornado público o conjunto de imagens sem o qual este trabalho não existiria.

Ao Instituto Federal do Espírito Santo, Campus Cachoeiro de Itapemirim, e aos professores do curso de Sistemas de Informação, pela formação.

Ao Sesi Cachoeiro e aos professores que me incentivaram a escolher este curso e com quem aprendi muito.

Aos colegas próximos do Ifes, pela convivência e pela companhia ao longo do curso.
